% Propietario recuperado íntegro: observador, automedida y Stokes holográfico finito.
% Notación propia del propietario original, preservada localmente para su consumo modular.
\providecommand{\X}{\mathcal X}
\providecommand{\Acal}{\mathfrak a}
\providecommand{\bd}{\partial}
\providecommand{\oneE}{\mathbf 1_E}

\section{Fundamentos algebraicos del cosido de fase y de la automedida}

\begin{definition}[Geometría de cosido de fase]
Una \emph{geometría de cosido de fase} es un sistema
\[
\X=(X,\oplus,\otimes,\tau,C_n,L,\mathfrak c,v),
\]
donde (X) es un soporte local finito; \(\oplus\) y \(\otimes\) son sus
dos leyes internas de transporte; \(\tau:X\to D\) es la aplicación residual;
\(C_n\) es un ciclo orientado; \(L\) es un retículo cuyo espacio real
\(W=L\otimes_{\mathbb Z}\mathbb R\) lleva un producto interno positivo;
\(\mathfrak c\in\mathbb P(D^n)\) es una clase residual distinguida no
trivial; y \(v\in L\) es su representante primitivo en la cámara fijada.
Un transporte \(\omega\in C^1(C_n;W)\) es \emph{admisible} cuando el modo
armónico de su descomposición pertenece a la recta
\(\mathbb Rv\subset W\).
\end{definition}

\begin{theorem}[Descomposición exacta--armónica sobre un ciclo]
\label{thm:hodge-cycle}
Para toda (1)-cocadena \(\omega\in C^1(C_n;W)\) existe una descomposición
única
\[
  \omega=d\Phi+u\oneE,
\]
con \(\Phi\in C^0(C_n;W)\), normalizada por \(\Phi(0)=0\), y \(u\in W\).
\end{theorem}

\begin{proof}
Si \(e_0,\ldots,e_{n-1}\) son las aristas orientadas del ciclo, se define
\[
 u=\frac1n\sum_{j=0}^{n-1}\omega(e_j),
 \qquad y_j=\omega(e_j)-u.
\]
Entonces \(\sum_jy_j=0\). Las sumas parciales
\(\Phi(k)=\sum_{j=0}^{k-1}y_j\), con \(\Phi(0)=0\), determinan una
(0)-cocadena bien definida y verifican \(d\Phi(e_j)=y_j\), incluida la
arista final. Para la unicidad, de
\(d\Phi+u\oneE=d\Psi+w\oneE\) se obtiene, tras sumar sobre el ciclo,
\(n(w-u)=0\); por tanto \(u=w\), y la normalización elimina la constante
restante de \(\Phi-\Psi\).
\end{proof}

\begin{definition}[Defecto de costura]
El levantamiento continuo a trozos \(\Gamma_\omega:[0,1]\to W\) se fija por
\[
 \Gamma_\omega\!\left(\frac{j+1}{n}\right)
 -\Gamma_\omega\!\left(\frac jn\right)=\omega(e_j).
\]
Su defecto de costura es
\[
 \bd\Gamma_\omega
 =\Gamma_\omega(1)-\Gamma_\omega(0)
 =\sum_{j=0}^{n-1}\omega(e_j).
\]
\end{definition}

\begin{proposition}[Anulación telescópica del término exacto]
\label{prop:exact-vanish}
Para toda (0)-cocadena \(\Phi\) se cumple
\[
 \sum_{j=0}^{n-1}(d\Phi)(e_j)=0.
\]
En consecuencia, el defecto de costura depende exclusivamente del modo
armónico \(u\oneE\) de la descomposición del
\cref{thm:hodge-cycle}.
\end{proposition}

\begin{proof}
La suma de \(\Phi(j+1)-\Phi(j)\) telescopa sobre el ciclo. Así, si
\(\omega=d\Phi+u\oneE\), entonces
\(\bd\Gamma_\omega=nu\).
\end{proof}

\begin{theorem}[Teorema universal de cosido]
\label{thm:universal-sewing}
Sea \(\X\) una geometría de cosido de fase y sea \(\omega\) un transporte
admisible. Existe un único escalar real \(\Acal_\X(\omega)\) tal que
\[
 \bd\Gamma_\omega=\Acal_\X(\omega)v,
 \qquad
 \Acal_\X(\omega)
 =\frac{\langle\bd\Gamma_\omega,v\rangle}{\langle v,v\rangle}.
\]
\end{theorem}

\begin{proof}
Por el \cref{thm:hodge-cycle}, \(\omega=d\Phi+u\oneE\), con \(u\) único.
La admisibilidad proporciona un único \(\lambda\in\mathbb R\) con
\(u=\lambda v\). La \cref{prop:exact-vanish} da
\(\bd\Gamma_\omega=n\lambda v\); basta definir
\(\Acal_\X(\omega)=n\lambda\) y proyectar ortogonalmente sobre
\(\mathbb Rv\).
\end{proof}

\section{Estado conjunto observador--observable, automedida intrínseca e identidad holográfica de Stokes}\label{sec:observer-holo}

Esta sección hace explícitas tres afirmaciones que en el resto del manuscrito aparecían
sólo en forma narrativa. La primera es que el \emph{observador} no introduce el invariante,
sino que selecciona una polarización, una trivialización y una orientación de lectura.
La segunda es que la \emph{auto-medida} no es una metáfora, sino el único funcional lineal
normalizado compatible con el modo residual de borde. La tercera es que el llamado
\emph{principio holográfico} admite, en este marco finito, una formulación exacta mediante
Stokes discreto: el mismo escalar puede leerse en el borde como defecto de costura o en el
\emph{bulk} como curvatura total proyectada.

Desde el punto de vista histórico, estas tres piezas permiten reescribir la noción clásica de
medida. En metrología estándar, una medida es trazable cuando puede relacionarse con una referencia
mediante una cadena documentada e ininterrumpida de comparaciones; y, desde 2019, el SI se define
por valores fijados de constantes de la naturaleza \cite{NISTTraceability2023,BIPMSI2019}.
El marco presente no niega ese esquema: lo desplaza hacia un nivel anterior, en el que la
"unidad" ya no es primero patrón externo, sino modo residual interno. La medida deja entonces de
ser comparación con un artefacto y pasa a ser lectura de un cierre propio.

\subsection{El observador como polarización, trivialización y orientación}

La interpretación geométrica más estable del problema del observador es la siguiente. El soporte
local posee al menos dos polarizaciones internas, una aditiva y otra multiplicativa. En el lenguaje
estándar de la cinemática cuántica finita, ambas polarizaciones generan una extensión de
Weyl--Heisenberg en la que la fase central registra el orden de los transportes
\cite{Tolar2014,Vourdas2006}. El observador no añade nueva dinámica al sistema; selecciona un
representante de la misma clase de transporte.

\begin{definition}[Observador de borde]
Sea $\X=(X,\oplus,\otimes,\tau,C_n,L,\mathfrak c,v)$ una geometría de cosido de fase y sea
$\omega\in C^1(C_n;W)$ un transporte admisible. Un \emph{observador de borde} es un par
\[
\mathcal O=(\varepsilon,\Phi),
\qquad
\varepsilon\in\{+1,-1\},
\qquad
\Phi\in C^0(C_n;W),
\]
que actúa sobre $\omega$ por
\[
\omega^{\mathcal O}:=\varepsilon\,\omega+d\Phi.
\]
El signo $\varepsilon$ codifica la orientación de lectura y la $0$-cochain $\Phi$ codifica la
trivialización o elección de gauge.
\end{definition}

\begin{remark}
En el nivel local, la elección entre lector aditivo y lector multiplicativo debe entenderse como una
polarización del mismo soporte de fase. En el nivel global de borde, esa elección queda absorbida en
la clase de la $1$-cochain observada. La parte geométricamente relevante es, por tanto, la clase de
cohomología y no el representante particular.
\end{remark}

\begin{proposition}[Covarianza del observador]\label{prop:observer-covariance}
Para todo transporte admisible $\omega$ y todo observador de borde
$\mathcal O=(\varepsilon,\Phi)$ se tiene
\[
\Acal_\X(\omega^{\mathcal O})=\varepsilon\,\Acal_\X(\omega).
\]
En particular, la auto-medida es invariante bajo cambio de trivialización y sólo cambia de signo al
invertir la orientación.
\end{proposition}

\begin{proof}
Usando la definición de defecto de costura y la linealidad, se obtiene
\[
\bd\Gamma_{\omega^{\mathcal O}}
=
\sum_{e\in C_n}\bigl(\varepsilon\,\omega(e)+(d\Phi)(e)\bigr)
=
\varepsilon\sum_{e\in C_n}\omega(e)+\sum_{e\in C_n}(d\Phi)(e).
\]
La suma del término exacto telescopa y vale $0$ por la
Proposición~\ref{prop:exact-vanish}. Luego
\[
\bd\Gamma_{\omega^{\mathcal O}}=\varepsilon\,\bd\Gamma_\omega.
\]
Proyectando sobre la recta residual $\R v$ mediante la fórmula del
Teorema~\ref{thm:universal-sewing}, se concluye
\[
\Acal_\X(\omega^{\mathcal O})
=
\frac{\langle \bd\Gamma_{\omega^{\mathcal O}},v\rangle}{\langle v,v\rangle}
=
\varepsilon\,\frac{\langle \bd\Gamma_\omega,v\rangle}{\langle v,v\rangle}
=
\varepsilon\,\Acal_\X(\omega).
\]
\end{proof}

\begin{corollary}[El observador no crea el invariante]\label{cor:observer-does-not-create}
Dos observadores de borde que difieran sólo por gauge miden el mismo escalar de cierre. Por tanto,
el observador fija una lectura, pero no fabrica el modo residual.
\end{corollary}

\begin{remark}
Esta proposición resuelve el problema del observador en el nivel exacto en que el presente marco
puede resolverlo. El observador no desaparece, pero su papel cambia: no introduce el patrón, sino
que selecciona una polarización y una orientación para un invariante ya generado por la propia
estructura. En este sentido, la teoría es \emph{covariante respecto del observador}.
\end{remark}

\subsection{Unicidad del funcional intrínseco de automedida}

El paso siguiente consiste en formalizar la idea de que el sistema se mide a sí mismo. Para ello hay
que distinguir entre la densidad local del modo armónico y la magnitud global de cierre del borde.

\begin{definition}[Densidad de auto-medida y cierre global]
Sea $\omega$ un transporte admisible sobre $C_n$. Definimos la \emph{densidad de auto-medida}
por
\[
\mu_\X(\omega):=\frac{1}{n}\,\Acal_\X(\omega)
=\frac{1}{n}\,\frac{\langle \bd\Gamma_\omega,v\rangle}{\langle v,v\rangle}.
\]
Equivalentemente, si
\[
\omega=d\Phi+u\oneE
\qquad\text{con }u\in\R v,
\]
entonces $u=\mu_\X(\omega)\,v$.
Llamaremos \emph{cierre global} al escalar
\[
\Acal_\X(\omega)=n\,\mu_\X(\omega).
\]
\end{definition}

\begin{lemma}[Cociente unidimensional de transportes admisibles]\label{lem:one-dimensional-quotient}
Sea $\mathcal T_{\mathrm{adm}}\subset C^1(C_n;W)$ el subespacio de transportes admisibles. Entonces
el cociente
\[
\mathcal T_{\mathrm{adm}}/dC^0(C_n;W)
\]
es un espacio vectorial real de dimensión $1$, canónicamente generado por la clase de
$\,v\oneE$.
\end{lemma}

\begin{proof}
Por el Teorema~\ref{thm:hodge-cycle}, todo $\omega\in\mathcal T_{\mathrm{adm}}$ se escribe de forma
única como
\[
\omega=d\Phi+u\oneE.
\]
La admisibilidad impone $u\in\R v$, de modo que la clase de $\omega$ módulo exactos queda
completamente determinada por el escalar real $\lambda$ tal que $u=\lambda v$. En consecuencia,
el cociente es isomorfo a $\R v$, y por tanto unidimensional.
\end{proof}

\begin{theorem}[Unicidad de la auto-medida]\label{thm:uniqueness-self-measurement}
Sea $F:\mathcal T_{\mathrm{adm}}\to\R$ un funcional lineal tal que:
\begin{enumerate}[label=\textup{(\roman*)}]
\item $F(\omega+d\Phi)=F(\omega)$ para toda $\Phi\in C^0(C_n;W)$;
\item $F(v\oneE)=1$.
\end{enumerate}
Entonces
\[
F=\mu_\X.
\]
En particular, la auto-medida es el \emph{único} funcional lineal normalizado compatible con el modo
residual distinguido.
\end{theorem}

\begin{proof}
Por el Lema~\ref{lem:one-dimensional-quotient}, el cociente
$\mathcal T_{\mathrm{adm}}/dC^0(C_n;W)$ es unidimensional y está generado por la clase de
$v\oneE$. Toda aplicación lineal invariante por gauge factoriza por dicho cociente y queda
completamente determinada por su valor sobre el generador. La normalización
$F(v\oneE)=1$ fija entonces unívocamente el funcional. Por definición,
\[
\mu_\X(v\oneE)=\frac1n\frac{\langle nv,v\rangle}{\langle v,v\rangle}=1,
\]
y además $\mu_\X$ es lineal e invariante por gauge. Luego $F=\mu_\X$.
\end{proof}

\begin{corollary}[Auto-medida en el caso HMT]\label{cor:self-measurement-hmt}
En la especialización dodecafásica fuerte, donde $n=12$ y $v=v_\alpha$, se tiene
\[
\mu_{\mathrm{HMT}}(\omega)=\frac{\alpha}{12},
\qquad
\Acal_{\mathrm{HMT}}(\omega)=\alpha.
\]
Así, el número $\alpha$ es el cierre global del borde y $\alpha/12$ es la densidad elemental de
auto-medida por fase.
\end{corollary}

\begin{remark}
Esta formulación cambia el significado de ``unidad''. En la metrología clásica, la unidad es una
referencia externa a la que un resultado se vincula mediante una cadena de calibraciones
\cite{NISTTraceability2023,BIPMSI2019}. Aquí la normalización proviene del propio sistema: el
vector $v$ no se importa desde fuera, sino que se genera internamente como lift mínimo de la clase
residual distinguida. En ese sentido preciso, la unidad es una \emph{auto-medida}.
\end{remark}

\subsection{Principio holográfico finito: borde, \emph{bulk} y Stokes discreto}

Lo que en el prólogo aparecía como intuición holográfica admite ahora una forma exacta. El hecho
crucial es que el mismo escalar puede leerse de dos maneras: como holonomía de borde o como
curvatura total proyectada del \emph{bulk}. Esto no identifica la teoría con la holografía
cuántico-gravitatoria de 't Hooft o Susskind \cite{tHooft1993,Susskind1995}; sí produce, en cambio,
un análogo finito y literal de reducción al borde.

\begin{definition}[\emph{Bulk} finito con borde cíclico]
Un \emph{bulk} finito con borde cíclico es un complejo celular orientado de dimensión $2$
\[
K=K^{(0)}\cup K^{(1)}\cup K^{(2)}
\]
cuya frontera orientada satisface
\[
\partial K\cong C_n.
\]
Dada una $1$-cochain $\Omega\in C^1(K;W)$, llamaremos \emph{curvatura discreta} a la
$2$-cochain
\[
F_\Omega:=d\Omega\in C^2(K;W).
\]
Si $\Omega\vert_{\partial K}=\omega$, diremos que $\Omega$ es una \emph{extensión de bulk} del
transporte de borde $\omega$.
\end{definition}

\begin{theorem}[Stokes discreto proyectado]\label{thm:discrete-stokes-projected}
Sea $K$ un bulk finito con frontera $\partial K\cong C_n$, y sea $\Omega\in C^1(K;W)$ una
extensión de bulk de un transporte admisible $\omega$ sobre la frontera. Entonces
\[
\bd\Gamma_\omega
=
\sum_{e\subset \partial K}\Omega(e)
=
\sum_{f\in K^{(2)}}F_\Omega(f).
\]
En particular,
\[
\Acal_\X(\omega)
=
\frac{\left\langle \sum_{f\in K^{(2)}}F_\Omega(f),v\right\rangle}{\langle v,v\rangle}.
\]
\end{theorem}

\begin{proof}
La primera igualdad es simplemente la definición de defecto de costura aplicada a la restricción de
$\Omega$ a la frontera. La segunda es la fórmula de Stokes discreto para cochains: la suma orientada
sobre la frontera de un complejo coincide con la suma, sobre todas las $2$-celdas, del coborde
orientado de la $1$-cochain. Es decir,
\[
\sum_{e\subset \partial K}\Omega(e)=\sum_{f\in K^{(2)}}(d\Omega)(f)=\sum_{f\in K^{(2)}}F_\Omega(f).
\]
Proyectando sobre la recta residual $\R v$ y usando el
Teorema~\ref{thm:universal-sewing}, se obtiene la fórmula final.
\end{proof}

\begin{corollary}[Obstrucción de \emph{bulk} plano]\label{cor:no-flat-bulk}
Si un transporte de borde admisible $\omega$ admite una extensión de bulk plana,
esto es, una $\Omega$ con $F_\Omega=0$, entonces
\[
\Acal_\X(\omega)=0.
\]
Equivalentemente, si $\Acal_\X(\omega)\neq 0$, toda extensión de bulk del borde debe portar
curvatura discreta no nula.
\end{corollary}

\begin{proof}
Es una consecuencia inmediata del Teorema~\ref{thm:discrete-stokes-projected}. Si
$F_\Omega=0$, entonces la suma proyectada de curvaturas vale $0$, y por tanto
$\Acal_\X(\omega)=0$.
\end{proof}

\begin{corollary}[Lectura holográfica finita]\label{cor:finite-holography}
El invariante de cierre puede leerse indistintamente como observable de borde o como curvatura
integrada del \emph{bulk}:
\[
\Acal_\X(\omega)
=
\frac{\langle \bd\Gamma_\omega,v\rangle}{\langle v,v\rangle}
=
\frac{\left\langle \sum_{f\in K^{(2)}}F_\Omega(f),v\right\rangle}{\langle v,v\rangle}.
\]
En este sentido preciso, el sistema satisface un principio holográfico finito.
\end{corollary}

\begin{remark}
El corolario anterior es la forma más literal y más fuerte del principio holográfico que el manuscrito
puede sostener con pruebas internas. No afirma una dualidad gravitatoria completa. Afirma algo más
modesto y exacto: el escalar físico-matemático relevante del sistema \emph{factoriza por el borde} y,
si se extiende a un \emph{bulk}, coincide con la curvatura total proyectada de ese \emph{bulk}.
Cuando $\Acal_\X(\omega)=\alpha$, la igualdad dice que $\alpha$ es simultáneamente una holonomía
de borde y una curvatura integrada.
\end{remark}

\begin{corollary}[Consecuencia geométrica para HMT]\label{cor:hmt-bulk-curvature}
En el caso HMT fuerte, todo relleno de bulk del ciclo dodecafásico que extienda el transporte de
borde canónico debe satisfacer
\[
\left\langle \sum_{f\in K^{(2)}}F_\Omega(f),v_\alpha\right\rangle
=
\alpha\,\langle v_\alpha,v_\alpha\rangle
=
54\alpha.
\]
En particular, el borde HMT no admite relleno plano; su \emph{bulk} porta necesariamente curvatura
total proyectada no nula.
\end{corollary}

\begin{proof}
Basta aplicar el Corolario~\ref{cor:finite-holography} con $v=v_\alpha$ y usar que
$\langle v_\alpha,v_\alpha\rangle=54$.
\end{proof}



\begin{remark}[Sobre la geometría del \emph{bulk}]
El Teorema~\ref{thm:discrete-stokes-projected} no identifica por sí solo la geometría métrica del
\emph{bulk}; sí impone una condición negativa muy fuerte: cuando $\Acal_\X(\omega)\neq 0$, el
relleno no puede ser exacta o discretamente plano. Si, además, se busca una completación global
compatible con la lectura modular del borde, el candidato natural deja de ser el disco euclídeo y pasa
a ser un \emph{bulk} de tipo modular. En la teoría clásica, las funciones modulares viven en el
semiplano superior $\mathbb H$ y el invariante de Klein $J(\tau)$ clasifica el toro complejo global
\cite{DLMFModular}. El resultado presente no demuestra todavía esa identificación, pero sí justifica
por qué la intuición de un \emph{bulk} no euclídeo es matemáticamente razonable.
\end{remark}

\subsection{Consecuencias estructurales de la polarización, la automedida y la identidad discreta de Stokes}

Las tres piezas anteriores permiten fijar, con un grado de precisión que antes sólo se esbozaba,
la interpretación de observador, auto-medida y holografía.

\begin{enumerate}[label=\textup{(\arabic*)}]
\item El \textbf{observador} es una elección de polarización, orientación y gauge. Cambia la lectura,
pero no el modo residual, salvo por el signo asociado a la orientación.
\item La \textbf{auto-medida} es el único funcional lineal normalizado sobre el espacio admisible de
transportes. No necesita patrón externo porque su normalización proviene del lift mínimo interno.
\item El \textbf{principio holográfico} adquiere una forma exacta: el escalar relevante puede leerse
íntegramente en el borde o íntegramente como curvatura total proyectada del bulk.
\end{enumerate}

Esta triple identificación permite reformular la tesis central del manuscrito de manera más fuerte:
no sólo $\alpha$ es un invariante archimedeano de borde, sino que, dentro del marco finito
considerado aquí, $\alpha$ es el primer escalar que una célula autoinercial de fase produce al
medirse a sí misma y al proyectar toda su curvatura efectiva sobre el modo residual distinguido.
